\section{Related Work}

The landscape of audio synthesis and restoration has evolved through distinct paradigms, ranging from manual feature engineering to fully automated deep learning pipelines. Early approaches in this domain relied heavily on vocoders that incorporated strong prior knowledge of signal processing but often suffered from limited expressivity \citep{engel2020ddsp}. In contrast, modern generative models typically operate directly in the time or frequency domains; however, these representations are inefficient as they fail to leverage existing perceptual and physical constraints inherent to sound generation. To address this gap, differentiable digital signal processing (DDSP) frameworks have emerged, aiming to bridge the divide between data-driven learning and domain-specific knowledge \citep{engel2020ddsp}. Concurrently, unsupervised restoration techniques have gained traction by applying statistical reasoning to reconstruct clean signals from corrupted observations without requiring paired training data or explicit likelihood models of corruption \citep{lehtinen2018n2n}. This shift towards learning mappings solely from degraded examples represents a significant departure from traditional supervised paradigms, enabling robust performance even when ground-truth references are unavailable.

A critical bottleneck in deploying these advanced architectures lies the sensitivity of neural networks to empirical choices regarding model structure and hyperparameters \citep{jaderberg2017pbt}. Traditional manual tuning is time-consuming and prone to error, prompting a surge in interest for automated design methodologies. Neural Architecture Search (NAS) has become a cornerstone technique for overcoming these limitations by automating the feature engineering process through end-to-end learning of hierarchical extractors \citep{elsken2019nas}. NAS methods generally fall into two categories: reinforcement learning-based approaches that evolve discrete architectures and differentiable relaxation techniques like DARTS, which formulate search as a continuous optimization problem solvable via gradient descent. While these methods have achieved remarkable success in image classification and language modeling \citep{liu2019darts}, their application to audio restoration remains an emerging frontier where the high-dimensional nature of waveform data presents unique scalability challenges.

Hyperparameter optimization (HPO) further complicates the training landscape, as performance is heavily dependent on choices such as loss functions and optimizer settings. Population Based Training (PBT) offers a pragmatic solution by asynchronously optimizing both model parameters and hyperparameters within a fixed computational budget \citep{jaderberg2017pbt}. This approach effectively utilizes parallel resources to jointly evolve the population of models, often outperforming static configurations found through grid search or random sampling. Furthermore, meta-learning strategies provide mechanisms for fast adaptation across diverse tasks with limited data samples \citep{finn2017maml}, a capability highly relevant when training on sparse datasets typical in audio restoration scenarios where clean reference signals are scarce.

Despite the breadth of literature covering NAS variants and HPO techniques, many studies focus narrowly on specific domains like computer vision or natural language processing without addressing the nuances of acoustic signal reconstruction \citep{elsken2019nas}. Additionally, while meta-learning frameworks offer theoretical advantages for few-shot learning, their integration with differentiable DSP components remains underexplored in practical audio applications. Our work distinguishes itself by synthesizing these disparate threads—combining unsupervised restoration principles with automated architecture search and hyperparameter tuning specifically tailored to the constraints of real-time audio synthesis. Unlike prior efforts that treat NAS as a black-box procedure, we explicitly incorporate domain knowledge into the search space definition, ensuring that discovered architectures respect physical signal properties while maximizing perceptual quality metrics such as residual reduction relative to ideal signals.
